\subsection{DERIVATIVE OF AN INVERSE FUNCTION}

PROPOSITION 6. Let \(f\) be a homeomorphism of an interval \(I \subset \mathbf{R}\) onto an interval \(J = f(I) \subset \mathbf{R}\) , and let \(g\) be the inverse homeomorphism \(^3\) . If \(f\) is differentiable at the point \(x_0 \in I\) , and if \(f'(x_0) \neq 0\) , then \(g\) has a derivative equal to \(1/f'(x_0)\) at \(y_0 = f(x_0)\) .

For each \(y \in J\) we have \(g(y) \in I\) and \(u = f(g(y))\) ; we thus can write \(\frac{g(y) - g(y_0)}{y - y_0} = \frac{g(y) - x_0}{f(g(y)) - f(x_0)}\) , for \(y \neq y_0\) . When y tends to \(y_0\) while remaining in J and \(\neq y_0\) , then \(g(y)\) tends to \(x_0\) remaining in I and \(\neq x_0\) , and the right-hand side in the preceding formula thus has limit \(1/f'(x_0)\) , since by hypothesis \(f'(x_0) \neq 0\) .

COROLLARY. If \(f\) is differentiable on \(\mathbf{I}\) and if \(f'(x) \neq 0\) on \(\mathbf{I}\) , then \(g\) is differentiable on \(\mathbf{J}\) and its derivative at each point \(y \in \mathbf{J}\) is \(1 / f'(g(y))\) .

For example, for each integer \(n > 0\) , the function \(x^{1 / n}\) is a homeomorphism of \(\mathbf{R}_+\) onto itself, is the inverse of \(x^n\) , and has derivative \(\frac{1}{n} x^{\frac{1}{n} - 1}\) at each \(x > 0\) .

One deduces easily, from prop. 5, that for every rational number \(r = p / q > 0\) the function \(x^{r} = (x^{1 / q})^{p}\) has derivative \(rx^{r - 1}\) at every \(x > 0\) .

Remarks. 1) All the preceding propositions, stated for functions differentiable at a point \(x_0\) , immediately yield propositions for functions which are right (resp. left) differentiable at \(x_0\) , when, instead of the functions themselves, one considers their restrictions to the intersection of their intervals of definition with the interval \([x_0, +\infty[\) (resp. \(] - \infty, x_0]\) ); we leave it to the reader to state them.

\begin{enumerate}
\def\labelenumi{\arabic{enumi})}
\setcounter{enumi}{1}
\tightlist
\item
  The preceding definitions and propositions (except for those concerning right and left derivatives) extend easily to the case where one replaces \(\mathbf{R}\) by an arbitrary commutative non-discrete topological field \(\mathbf{K}\) , and the topological vector spaces (resp. topological algebras) over \(\mathbf{R}\) by topological vector spaces (resp. topological algebras) over \(\mathbf{K}\) . In def. 1 and props. 1, 2 and 3 it is enough to replace I by a neighbourhood of \(x_0\) in \(\mathbf{K}\) ; in prop. 4 one must assume further that the map \(\mathbf{y} \mapsto \mathbf{y}^{-1}\) is defined and continuous on a neighbourhood of \(\mathbf{f}(x_0)\) in E. Prop. 5 generalizes in the following manner: let \(\mathbf{K}'\) be a non-discrete subfield of the topological field \(\mathbf{K}\) , let E be a topological vector space over \(\mathbf{K}\) ; let \(f\) be a function defined on a neighbourhood \(V \subset \mathbf{K}'\) of \(x_0 \in \mathbf{K}'\) , with values in \(\mathbf{K}\) (considered as a topological vector space over \(\mathbf{K}'\) ), differentiable at \(x_0\) , and let \(\mathbf{g}\) be a function defined on a neighbourhood of \(f(x_0) \in \mathbf{K}\) , with values in E, and differentiable at the point \(f(x_0)\) ; then the map \(\mathbf{g} \circ f\) is differentiable at \(x_0\) and has derivative \(\mathbf{g}'(f(x_0))f'(x_0)\) there (E being then considered as a topological vector space over \(\mathbf{K}'\) ).
\end{enumerate}

With the same notation, let \(\mathbf{f}\) be a function defined on a neighbourhood \(V\) of \(a \in K\) , with values in \(E\) , and differentiable at the point \(a\) ; if \(a \in K'\) , then the restriction of \(\mathbf{f}\) to \(V \cap K'\) is differentiable at \(a\) , and has derivative \(\mathbf{f}'(a)\) there. These considerations apply above all, in practice, to the case where \(K = C\) and \(K' = R\) .

Finally, prop. 6 extends to the case where one replaces I by a neighbourhood of \(x_0 \in \mathbf{K}\) , and \(f\) by a homeomorphism of I onto a neighbourhood \(\mathrm{J} = f(\mathrm{I})\) of \(y_0 = f(x_0)\) in \(\mathbf{K}\) .

